\section{Cartesian control}
In Cartesian control, the task is to control a frame
associated with the robot to either reach or track a desired frame.
For manipulators, the frame of interest is the end-effector frame 
$\left\{e\right\}$,
as it has the tool (usually a gripper) which the robot uses
to interact with the environment. If we want the robot to pick
up an object, the target frame $\left\{goal\right\}$ 
is usually defined in the center
of the object, with an appropriate orientation for grasping.
For mobile bases, the frame is usually defined in the center of
the mobile base. 

The problem of determining the pose $T_w^e$ 
(end-effector frame location
as seen in the world frame) of the end-effector
from the angle of joint positions $q$ is referred
to as \textbf{forward kinematics},
$T_w^e = f(q)$.
However, we usually care about the inverse: knowing which 
joint angles $q$ correspond to an end-effector pose $T_w^e$.
This problem is referred to as \textbf{inverse kinematics}.
The issue is that the inverse of the forward kinematics map $f^{-1}$
is ill-defined unless the robot has 6 degrees of freedom. 
More specifically, frames, poses and transformations all life
in a 6-dimensional space $\mathbb{SE}(3)$.
Without getting lost in the math, $\mathbb{SE}(3)$ is not a vector space,
$\mathbb{SE}(3) \neq \mathbb{R}^6$, but it's still 6-dimensional.
For example a sphere can be described with 2 coordinates (lattitude and 
longitude), so it's 2-dimensional, despite being embedded in 3D-space.
Point being, if we have a robot with more than 6 degrees of freedom,
and the task is defined in 6 dimensions, we have an underdetermined
problem (the set of equations is smaller than the number of unknowns).

Since we can't use $f^{-1}$ directly, we need a different strategy,
which is referred to as \textbf{differential inverse kinematics}.

\paragraph{Differential inverse kinematics}
Instead of working with the forward kinematic map,
and its inverse, we work its derivative. This allows 
us to solve the differential inverse kinematics problem,
which is to find the joint velocities $\dot{q}$ such that
the end-effector has velocity $\mathcal{V}$.
We can then re-solve this problem at every time instant,
and use the computed joint velocities $\dot{q}$ to guide the 
end-effector to the goal.

In math it goes as follows. First we take the derivative
of the forward kinematics function:
\begin{gather}
		T_{ w }^{ e } = f (q) \quad \Big/ \frac{d }{d t} \\
\mathcal{V} = J (q) \dot{q}
\end{gather}
The Jacobian (differential) of the forward kinematics map
$ J = \frac{d f}{dq}   $
is a crucial object which will come up in many contexts.
But for the task at hand, it is critical to notice
its dimensions:
$ J \in \mathbb{R}^{6 \times nq}  $ where $ nq  $ is the number of joints.
Whereas $ \mathcal{V} \in \mathbb{R}^{6}  $.
Thus, in general, the Jacobian $ J  $ is not a square matrix,
and thus has no defined inverse.
How this is dealt with is what defines different
differential inverse kinematics solvers. The provided
options are explored in \ref{sub-ik-solvers}.
But the simplest one is use a generalization of matrix
inverses to non-square matrices called the pseudoinverse (Moore-Penrose
inverse)
$ J^{ \dagger } = J^{ T } (JJ^{ T })^{ -1 }  $.
This leads us to the simplest differential inverse kinematics solution
\begin{equation}
		\dot{q} = J^{ \dagger } \mathcal{V}
\end{equation}
Equipped with this, we can first compute the velocity
the end-effector needs to get to a desired location,
solve for joint velocities and command them.
Doing so until the robot converges to the target location
yields a controller solving inverse kinematics,
which is explored next.

\subsection{Cartesian space velocity control}
\subsubsection{Point-to-point motion}
In point-to-point motion, we compute difference (error) between
the current end-effector position $ T_{ w }^{ e }  $ 
and the target position $ T_{ w }^{ goal }  $. We then set a velocity
parallel to this direction as the desired end-effector velocity,
and solve the differential inverse kinematics problem.
The error vector is given by
\begin{equation}
e = \log_{ 6 } \left( \left(   T_{ w }^{ e }\right)^{ -1 } \cdot T_{ w }^{ \text{goal} }  \right) 
\end{equation}
which means $ e  $ is a twist (6D velocity), i.e. $ e \in se (3)  $.
If you want to understand what this expression means, you have
to become familiar with Lie groups. The crash course version
covering everything you need 
is given in \ref{sec-robotics-lie-algebra}.

Having the desired twist, we call an \fbox{\texttt{ik\_solver}},
which is specified with an argument of the same name.
The available inverse kinematics solvers are covered in 
\ref{sub-ik-solvers}.
The algorithm goes as follows

\begin{algorithm}[H]
\caption{moveL}
\begin{algorithmic}[1]
		\Procedure{moveL}{cfg, robot, $ T_{ w }^{ goal }  $}
		\While{$\left\lVert   \mathcal{V}_{ cmd }\right \rVert >$ goal\_error }
		\State $   E \gets \left( \left(T_{ w }^{ e }\right)^{ -1 } \cdot T_{ w }^{ \text{goal} }  \right) $
		\State $\mathcal{V}_{ cmd } \gets \log (E)$
		\State $v_{ cmd } \gets $ ik\_solver (cfg, robot, $ \mathcal{V}_{ cmd }  $)
		\EndWhile
		\EndProcedure
\end{algorithmic}
\end{algorithm}


\subsection{A deeper look into inverse kinematics}
\subsubsection{A key insight from linear algebra}
The Jacobian $J$, as any matrix, can be interpreted as a linear operator.
Since $J \in \mathbb{R}^{6 \times nq}$, it is a linear map 
$J: \mathbb{R}^{nq} \mapsto se(3) = \mathbb{R}^{6}$.
From linear algebra, we know that a map can ``cover'' only as many dimensions
as it has inputs. More formally, the dimension of its image can not be higher 
than the dimension of its domain.
What happens to the other input dimensions? Since they do not affect the image space,
they are all mapped to 0.
Critically, we can split the linear map into the operation which projects onto the image,
and the map which projects into 0. This has direct consequences to inverse kinematics
as the map which projects into 0 defines the set of joint velocities for which the end-effector does not move.
This is then leveraged for better control of the robot.
In order to see how this is done, and what this map splitting even means, we need to introduce
some notation and state a theorem.

Let:
\begin{itemize}
\item $L$ be some linear map from vector space $v$ to vector space $W$, $L: V \mapsto W$
\item $Im(L)$ be the image of $L$, i.e. the subspace of $W$ to which $L$ maps to, $Im(L) \subset W$
\item $\mathcal{N} (L)$ be the nullspace (also called kernel) of $L$, i.e. the subspace of $ L \subset V$ which $L$ maps to 0 $\mathcal{N}(L): V \mapsto 0 $
\item $dim()$ denote the dimension of a vector space
\end{itemize}

The rank nullity theorem states:
\begin{equation}
dim(Im(L)) + dim(\mathcal{N}(L)) = dim(V)
\end{equation}

\paragraph{What does this mean in the context of inverse kinematics?}
For a robot with $nq > 6$, the will be exactly $nq - 6$ additional degrees of freedom.
These additional degrees of freedom are referred to as \textbf{redundant},
and collectively \textbf{redundancy}.
For some robot pose $q$, the nullspace of its Jacobian $\mathcal{N}(J)$ defines the set of 
of joint velocities for which the end-effector does not move.
This can be used to put a robot in a more task-compatible pose (exert more force in a direction,
avoid an obstacle etc). Defining additional tasks which are meant to be achieved
in the Jacobian nullspace is called \textbf{redundancy resolution}.
Sometimes this can be just simple joint-limit avoidance, or trying to maintain
some desired posture $q_{ref}$.
The more redundancy a robot has, the more important it is to have appropriate redundancy resolution.

Pseudoinverse-based inverse kinematics algorithms directly construct secondary tasks for redundancy resolution
in the nullspace, which is covered in the paragraph on the pseudoinverse IK solver (TODO: write that down).
In other methods, for example in the QP formulation, this is done indirectly.


\subsubsection{IK as an optimization problem}
The inverse kinematics problem can be posed as an optimization problem.
In fact, there is a strong case to be made that turning an underdefined
set of equations $\mathcal{V} = J \dot{q}$ into an optimization problem is the only sensible thing to do.
The reason why is that a secondary task can be specified, leading to a full determined problem.
In fact, the pseudoinverse solution comes from using linear least-squares
to solve $min_{\dot{q}} \left\lVert \mathcal{V}\right \rVert$.

Here it is critical to observe that $min_{\dot{q}} \left\lVert \mathcal{V}\right \rVert$ is a quadratic cost function.
This means that inverse kinematics can be formulated as a \textbf{quadratic program} (QP).
Here's why this is so powerful: 
\begin{itemize}
\item QPs are convex, meaning we always find the true optimum, and can be solved incredibly quickly with dedicated solvers
\item QPs can handle linear equality and inequality constrains
\item a sum of quadratic costs is still a quadratic cost
\end{itemize}
At first glance you might not be impressed with linear constraints. But remember that we are doing differential
inverse kinematics, which is already a linearization of the inverse kinematics problem,
and that we solve it at every timestep. This means that we can get a lot of mileage out
of linear approximations of non-linear constraints.
Secondly, norms squared are quadratic. Thus we can define a QP for arbitrary obstacle avoidance.
While the Jacobian and the end-effector velocity define the main tasks, that does not mean
that other tasks need anything to do with either --- anything depending on $\dot{q}$ is fair game.
Finally, there is a crucial trick to have in mind, which is that
$\dot{q} \Delta t$ gives us a difference is position,
and $\dot{q}/\Delta  t$ gives us acceleration. 
You can simply add the current joint values to the difference is position,
and this then opens up the possibility to define costs and constraints in terms
of joint positions (ex. a constraint on joint limits).
This is legitimate as we can ensure that the solution is a feasible velocity command.

Here it is important to note that the different objectives (costs) need not be a primary and 
a secondary tasks --- you can have competing primary tasks.
This is because one can transcribe a constraint into a cost function via
\textbf{barrier functions}. This is extremely powerful as it allows
for constraints other than linear ones to be brought into the problem.
More on this in TODO: citation (caron's tutorials and pink code).
The important takeaway is that IK as QP is extremely flexible and powerful.
Essential, they are trajectory optimization with an infinitesimally small time horizon.
And while having a time horizon is often beneficial, many seemingly complex
control tasks can be solved perfectly well with just a single QP.



\subsection{Inverse kinematics solvers}
\label{sub-ik-solvers}
The code for all inverse kinematics solvers
is in 
\newline
\fbox{\texttt{python/smc/control/cartesian\_space/ik\_solvers.py}}.

\subsubsection{Jacobian-based}
There are of course the classics:
\begin{itemize}
	\item Jacobian pseudoinverse: 
			$ v = J^{ \dagger } V_{ e }= J^{ T }(JJ^{ T }
			+ \lambda I)^{ -1 } 
			V_{ e }$ TODO citation
	\item Jacobian transpose: $ v = J^{ T } V_{ e }  $ TODO citation
\end{itemize}

For larger movement, these methods can result in wonky 
trajectories as the error twist $ V_{ e }  $ is large,
resulting in large joint velocity commands $ v  $.
Those in turn get saturated, leading to the end-effector
not going in a straight line, but something else.
This could be solved by scaling the error twist
$ V_{ e }  $ to a smaller number, but this is 
not an ideal solution because the maximum end-effector 
twist in a particular direction depends on the configuration
(which could be taken into account manually, but there are 
better methods out there).


%\paragraph{QP formulation}
%The generic QP formulation is
%\begin{align}
%		\min_{x} & \frac{1}{2} x^{ T } P x + r^{ T } x \\
%		\text{s.t. } & G x \leq h
%					 & A x = b\\
%					 & x_{ \min } \leq x \leq x_{ \max }
%\end{align}
%
%The current one is wrong because I didn't know any better in my youth
%\begin{align}
%		\min_{v} & v^{ T } I v + r^{ T }x \\
%		\text{s.t. } & J v = V
%		& v_{ \min } \leq v \leq v_{ \max }
%\end{align}
%where $ r  $ could by the derivative of the manipulatibility index
%$ \frac{\partial m}{\partial q}  $, or $ q_{ \text{comfy\_home} } - q  $.
%
%The reason why this doesn't work is because $ V  $ is big enough,
%the solution $ v  $ becomes infeasible due to the constraint on it.
%If you remove the constraint, you can get shitty results due to
%the saturation on $ V  $. If you introduce a scaling $ KV  $
%with $ K   $ you get a shitty hot-fix.

\subsubsection{QP formulation}
\label{qp-formulation}
\begin{align}
		\min_{v} &\frac{1}{2} \left\lVert J \Delta q - \mathcal{V} \right \rVert  \\
		\text{s.t. } & v_{ \min } \leq v \leq v_{ \max }
\end{align}
The joint displacement $ \Delta q  $ can be regarded as joint velocity $ v  $
if the integration step is small enough (it is, i.e. we make sure 
that we have high enough control frequency),
$v = \Delta q/\Delta t$.
From now on we write $ v  $.
If there were no additional objectives, and no any constraints whatsoever,
the linear least-squares solution to this problem would be the Jacobian pseudoinverse method.
What makes the QP formulation superior is the fact that it does handle constraints.

This can be transcribed as a generic quadratic program (QP), which looks like
\begin{align}
		\min_{x} & \frac{1}{2} x^{ T } P x + q^{ T }x \\
		\text{s.t. } & Gx \leq h\\
					 & Ax = b
\end{align}

In particular, the transcription goes as follows:
\begin{align}
		\left\lVert J v - K V \right \rVert ^{ 2 } &=
		\left( J v - K V \right) ^{ T } \left( J v - K V \right) \\
&=
		v^{ T }J^{ T }Jv - KV^{ T }Jv - KvJ^{ T }V + 
		\underbrace{  K^{ 2 }V^{ T }V}_{\text{does not depend on }v} \\
		 &\implies v^{ T } \underbrace{  (J^{ T }J)}_{P}v 
	 \underbrace{  - 2 KV^{ T }J}_{q} v
\end{align}
\textbf{Note} that the same Tikhonov regularization (also known as Levenberg-Marquardt damping
in this context) 
should added to the $ P  $ as it was in the pseudoinverse, i.e.
$ P = J^{ T }J + \lambda I $.
Crucially, we can add additional tasks and constraints to the QP formulation,
as long as the costs are quadratic, and the constraints are linear.
Since we are solving this problem at each control time step, 
this is actually quite powerful as we can use linearised nonlinear costs/constraints
and they will work fine.

In particular, let a cost be denoted as $ e (q)  $,
and it's Jacobian by $ J (q)  $.
This will again result in a cost of form
$ \left\lVert J_{ i } v - K_{ i }e_{ i } \right \rVert ^{ 2 }  $.
You then work out the same steps as above for the frame placement cost, getting $ (P_{ i }, r_{ i })  $. 
Then you can simply put $ P = \sum_{i}^{} P_{ i }, r = \sum_{i}^{} r_{ i }  $,
which is still quadratic.
This emulates behaviour of HQP.
But why do we even want secondary objectives? Tackled in the next subsection.


Available QP solvers:
\begin{itemize}
		\item solvers supported by \href{https://github.com/quadprog/quadprog}{quadprog} (we use what the default one is,
				but you can install and select any one supported
				by quadprog)
		\item \href{https://github.com/Simple-Robotics/proxsuite}{proxsuite} - cool because it has warmstarting
				(among other things)
\end{itemize}

\subsubsection{Secondary tasks and redundancy resolution}
TODO:
\begin{itemize}
		\item what is redundancy
		\item nullspace exploitation
		\item why is that necessary for good performance
		\item can you have multiple primary objectives tho??
		\item constraints vs objectives
\end{itemize}

%Finding an appropriate secondary task, and doing that computation is a part of
%the TODO. Since the biggest issue on YuMi is the elbows hitting the torso,
%going for $ e = q_{ \text{comfy\_configuration} } - q   $ is probably the way to go.
%The alternatives are manipulability maximization or self-collision avoidance.
%Manipulability maximization could also be added as a ternary goal,
%where the manipulability is defined for the whole of the robot.
%Looking at \href{ https://github.com/stephane-caron/pink}{pink} 
%and  \href{https://stephane-caron.github.io/pink/index.html}{it's documentation}
%is a good place to start.
%\textbf{Note} that pink isn't integrated into SMC simply because it was annoying
%to match it's structure to SMC's, but it can be done.
%A starting point for that approach would be to look into
%\fbox{\texttt{examples/cartesian\_space/yumi\_self\_collision.py}}.

%\paragraph{Solution approach}
%We need the QP formulation of inverse kinematics,
%(or something which allows for additional objectives and constraints), 
%because the pseudoinverse solver
%causes the configuration of the robot goes to shit (breaks the robot) during
%Cartesian path-following (and thus cart-pulling as well).
%
%
%One way to implement the QP formulation is to shove the whole of pink setup into an \fbox{\texttt{ik\_solver}}
%construction. Then it can be applied to any \fbox{\texttt{controlLoop}} without any additional code. 
%That requires creating a big struct for all pink's variables,
%which is then partialed to comply with \fbox{\texttt{ik\_solver}}'s API.
%The API could also be changed, just note that you then need to refactor this throughout the codebase.
%
%The way adding tasks is handled in pink is great, and there is no point
%in rewriting that. We want this because it would make it super
%easy to swap different objectives, i.e.
%we wouldn't need to a new ik solver for every thing to try
%(or write the logic equivalent to pink to make one of the customizable).
%
%Alternatively, we can manually construct the QP, and also partial in
%whatever additional variables we need. The drawback is that
%this will be more verbose and annoying, but the upside is that
%the math of the QP will be completely transparent.



%\paragraph{Secondary goals}
%\begin{itemize}
%		\item posture reference tasks with $ e (q) = q_{ \text{comfy} } - q  $ --> DONE
%				\newline
%				ik\_solver is called \fbox{\texttt{QPWithPosture}}.
%				It does not take mobile pases positions into account of course.
%		\item velocity limit constraints --> DONE
%		\item TODO: joint limit constraints 
%		\item TODO IF ABOVE ISN'T ENOUGH FOR GOOD POSTURE: put in the derivative of manipulability w.r.t. $ q  $.
%		\item TODO IF SELF-COLLISIONS ARE COMMON: adding collision avoidance via control barrier functions.
%				Regarding what to avoid, there are 2 options:
%				\begin{enumerate}
%						\item make spheres around elbows and end-effectors. Use \fbox{\texttt{BodySphericalBarrier}}
%								from pink for this. Inputs are frames that can't be avoided and the radius
%								of the sphere around said frames. Look at 
%								\newline
%								\fbox{\texttt{pink/examples/barriers/yumi\_self\_collision.py}} for details.
%								There's an SMC-ed version of this in examples/cartesian\_control/yumi\_self\_collision.py
%						\item make ellipses around every link. Use \fbox{\texttt{SelfCollisionBarrier}}
%								from pink for this. Look at 
%								\newline
%								\fbox{\texttt{pink/examples/barriers/kukas\_self\_collision.py}} for details.
%				\end{enumerate}
%\end{itemize}

\paragraph{How to run}
The main file for SingleArmInterface point to point Cartesian control is 
\fbox{\texttt{example/cartesian\_space/clik\_point\_to\_point.py}},
and the one for DualArmInterface is 
\fbox{\texttt{example/cartesian\_space/dual\_arm\_clik.py}}.
The relevant arguments can be found simple by adding
\fbox{\texttt{--help}} when running the example,
but the most important ones is
\fbox{\texttt{--ik-solver [solver]}} to select a solver.

\subsection{Path following Cartesian velocity control}
TODO:
first we need to discuss what we mean by plan here

\subsubsection{[TODO]: path interpolation}
TODO: reword, put in a lil on B-splines


%This is the key thing to do to make sure Cartesian path following works.
%The first thing to do is determine the number of points to be interpolated between.
%There is no benefit in interpolating the whole path as
%we're not following the whole path.
%We want a simple curve to make the computation quick, as it needs to happen
%in real-time in conjunction with the real-time path planner
%(anything reasonable should be fine).
%\textbf{Note} that linear interpolation (without splines) isn't good enough
%as it results in jagged motion on the ``edges''.
%
%
%One option is to form a Bezier curve, or a polynomial interpolation on the first few points
%and call it a day. A smarter version would be 
%to do the same, but with least-squares fit on a larger number of points,
%the amount of which is determined by some heuristic like their arclength
%being at least something.



\paragraph{Short note on 2D to 6D reference}
Often we get a 2D path to follow, which is of course underdefined.
This is most common when drawing a path for the purposes of testing a path-following controller.
The path is made to be 3D by specifying a height $ h  $ (this is an argument),
giving the translation $ p = \begin{bmatrix} x_{ i } & y_{ i } & h \end{bmatrix}   $
Of course, we could impose an arbitrary orientation along the path,
and interpolate the $ SO (3)  $ rotations to make it smooth. 
However, this isn't implemented as it (so far) isn't useful for cart-pulling nor navigation.

Instead, a Frenet frame is constructed on the linear path to give it orientation.
First, we make the x-axis be the tangent of the path. 
The angle of the x-axis of the frame as seen in the x-y plane of
the world frame $ \left\{ w \right\}   $ is determined with
\begin{equation}
 \theta [i] = atan2 \left(  \frac{\phi_{ y } [i+1] - \phi _{ y } [i]}{\phi_{ x } [i+1] - \phi _{ x } [i]}  \right)  
\end{equation}
and we turn this into a rotation matrix with
$ R_{ \theta_{ i } }  \text{rpyToMatrix} (0,0,\theta_{ i })$
Next, we specify the rotation about the x-axis (this is an argument),
and get $ R_{ x } =  \text{rpyToMatrix} (\text{rot\_x},0,0)$.
In cart-pulling, we want the z axis to be perpendicular to the 
The y-axis completes the right-handed frame, meaning
$ R = R_{ \theta } \cdot R_{ x }  $.
Finally, the resulting $ SE (3)  $ path is then
\begin{equation}
P_{ i } = 
\begin{bmatrix} 
		R_{ i } & p_{ i } \\
		0 & 1
\end{bmatrix} 
\end{equation}


\subsubsection{Trajectory generation}
TODO

\subsubsection{Velocity tracking controller}

\begin{algorithm}[H]
\caption{cartesianTrajectoryTracking}
\begin{algorithmic}[1]
		\Procedure{cartesianTrajectoryTracking}{cfg, robot, $ T_{ w }^{ ref } (t)  $}
		\While{$t < T_{ final } $ and $\left\lVert   \mathcal{V}_{ cmd }\right \rVert >$ goal\_error }
%		\State $   E \gets \left( \left(T_{ w }^{ e }\right)^{ -1 } \cdot T_{ w }^{ \text{goal} }  \right) $
%		\State $\mathcal{V}_{ cmd } \gets \log (E)$
%		\State $v_{ cmd } \gets $ ik\_solver (cfg, robot, $ \mathcal{V}_{ cmd }  $)
		\EndWhile
		\EndProcedure
\end{algorithmic}
\end{algorithm}


\paragraph{How to run}
The main file for SingleArmInterface path following Cartesian control is 
\fbox{\texttt{example/cartesian\_space/clik\_path\_following.py}},
and the one for DualArmInterface is 
\fbox{\texttt{example/cartesian\_space/dual\_arm\_clik\_path\_following.py}}.
The relevant arguments can be found simple by adding
\fbox{\texttt{--help}} when running the example,
but the most important ones is
\fbox{\texttt{--ik-solver [solver]}} to select a solver.
Additionally, there is a 
\fbox{\texttt{--planner}} option (on by default).
This will spawn the real-time path generation 
in a new process, and create a predefined map.

\subsection{Cartesian space torque control}
TODO
